\documentclass[12pt]{article}
\usepackage{amsmath}
\usepackage{amsfonts}
\usepackage{amssymb}
\usepackage{physics}
\usepackage{geometry}
\geometry{a4paper, margin=0.7in}


\title{Homework assignment for 02YMECH \\ Chapter: Mathematical tools}
\date{}
\begin{document}
\maketitle

\section*{Solutions}

\begin{enumerate}

% 1
\item \textbf{Indefinite integrals.}

\begin{enumerate}
    \item Compute
    \[
      \int \frac{dx}{a+x}.
    \]
    Let $u = a+x$, so that $du = dx$. Then
    \[
      \int \frac{dx}{a+x}
      = \int \frac{du}{u}
      = \ln|u| + C
      = \ln|a+x| + C.
    \]

    \item Compute
    \[
      \int \left(3x^2 e^{x^3} + \frac{4x}{1+x^2}\right) \, dx.
    \]

    We split the integral:
    \[
      \int 3x^2 e^{x^3} \, dx + \int \frac{4x}{1+x^2} \, dx.
    \]

    For the first term, set $u = x^3$, $du = 3x^2 \, dx$, so
    \[
      \int 3x^2 e^{x^3} \, dx = \int e^u \, du = e^u + C = e^{x^3} + C.
    \]

    For the second term, set $v = 1+x^2$, $dv = 2x \, dx$, so $4x \, dx = 2\, dv$. Then
    \[
      \int \frac{4x}{1+x^2} \, dx
      = \int \frac{2\, dv}{v}
      = 2 \ln|v| + C
      = 2 \ln(1+x^2) + C,
    \]
    since $1+x^2>0$.

    Therefore,
    \[
      \int \left(3x^2 e^{x^3} + \frac{4x}{1+x^2}\right) \, dx
      = e^{x^3} + 2 \ln(1+x^2) + C.
    \]

\end{enumerate}


% 2
\item \textbf{Definite integral.}

Compute
\[
  \int_{1}^{4} (2x+3)\,dx.
\]

An antiderivative of $2x+3$ is
\[
  \int (2x+3)\,dx = x^2 + 3x + C.
\]

Evaluate at the limits:
\begin{align*}
\int_{1}^{4} (2x+3)\,dx
&= \left[ x^2 + 3x \right]_{1}^{4} \\
&= \left(4^2 + 3\cdot 4\right) - \left(1^2 + 3\cdot 1\right) \\
&= (16+12) - (1+3) \\
&= 28 - 4 \\
&= 24.
\end{align*}

Interpretation: this is the area under the straight line $2x+3$ between $x=1$ and $x=4$.


% 3
\item \textbf{Derivative of a function.}

Given
\[
  f(x) = \frac{x^3 \ln(x)}{e^{2x}} = x^3 \ln(x)\, e^{-2x}, \qquad x>0.
\]

We differentiate using the product rule on $x^3 \ln(x)$ and $e^{-2x}$, or equivalently use the product rule plus chain rule.

First compute
\[
g(x) = x^3 \ln(x), \qquad h(x) = e^{-2x}, \qquad f(x)=g(x)h(x).
\]

Then
\[
  g'(x) = \dv{x} \left( x^3 \ln x \right)
  = 3x^2 \ln x + x^3 \cdot \frac{1}{x}
  = 3x^2 \ln x + x^2
  = x^2 (3\ln x + 1).
\]

Also
\[
  h'(x) = \dv{x}\left(e^{-2x}\right) = (-2)e^{-2x}.
\]

Product rule:
\begin{align*}
f'(x)
&= g'(x)h(x) + g(x)h'(x) \\
&= \left[x^2 (3\ln x + 1)\right] e^{-2x} + \left[x^3 \ln x\right] (-2 e^{-2x}) \\
&= e^{-2x} \left[ x^2 (3\ln x + 1) - 2 x^3 \ln x \right].
\end{align*}

Factor $x^2$:
\[
f'(x)
= x^2 e^{-2x} \left[ (3\ln x + 1) - 2x \ln x \right].
\]

So
\[
  f'(x) = x^2 e^{-2x} \left( 3\ln x + 1 - 2x \ln x \right), \qquad x>0.
\]

Explanation: we used product rule and the fact that $\dv{}{x}(\ln x)=1/x$ for $x>0$.


% 4
\item \textbf{Local extrema via first and second derivative tests.}

Let
\[
  f(x) = x^3 - 6x^2 + 9x + 1.
\]

First derivative:
\[
  f'(x) = 3x^2 -12x + 9.
\]

Set $f'(x)=0$:
\[
  3x^2 -12x + 9 = 0
  \quad\Longrightarrow\quad
  x^2 -4x +3=0
  \quad\Longrightarrow\quad
  (x-1)(x-3)=0.
\]
So critical points are $x=1$ and $x=3$.

Second derivative:
\[
  f''(x) = 6x -12.
\]

Evaluate $f''$:
\[
  f''(1)=6(1)-12 = -6 <0,
\]
which means $f$ is concave down at $x=1$, so $x=1$ is a local maximum.

\[
  f''(3)=6(3)-12 = 6>0,
\]
which means $f$ is concave up at $x=3$, so $x=3$ is a local minimum.

Now compute the function values:
\begin{align*}
f(1) &= 1^3 -6(1)^2 + 9(1) +1
= 1 -6 +9 +1
= 5, \\
f(3) &= 3^3 -6(3)^2 + 9(3) +1
= 27 -54 +27 +1
= 1.
\end{align*}

Conclusion:
\[
  \text{local maximum: } f(1)=5 \text{ at } x=1,
  \qquad
  \text{local minimum: } f(3)=1 \text{ at } x=3.
\]

Explanation: the sign of $f''$ at the critical point determines if the point is a max ($f''<0$) or a min ($f''>0$).


% 5
\item \textbf{Taylor series of $\sin x$, $\cos x$, $e^x$ about $0$ up to $x^5$.}

Maclaurin series (Taylor about $0$):

\[
\sin x = x - \frac{x^3}{3!} + \frac{x^5}{5!} + \mathcal{O}(x^7)
= x - \frac{x^3}{6} + \frac{x^5}{120} + \mathcal{O}(x^7).
\]

\[
\cos x = 1 - \frac{x^2}{2!} + \frac{x^4}{4!} + \mathcal{O}(x^6)
= 1 - \frac{x^2}{2} + \frac{x^4}{24} + \mathcal{O}(x^6).
\]

\[
e^x = 1 + x + \frac{x^2}{2!} + \frac{x^3}{3!} + \frac{x^4}{4!} + \frac{x^5}{5!} + \mathcal{O}(x^6)
= 1 + x + \frac{x^2}{2} + \frac{x^3}{6} + \frac{x^4}{24} + \frac{x^5}{120} + \mathcal{O}(x^6).
\]

Explanation: we truncated after the $x^5$ term or the highest power $\le 5$ that actually appears.


% 6
\item \textbf{Partial derivatives and total differential.}

Given
\[
  f(x,y) = x^2 y + e^{xy}.
\]

Partial derivative with respect to $x$:
\[
  f_x = \pdv{f}{x}
  = 2x y + e^{xy} \cdot \pdv{}{x}(xy)
  = 2xy + e^{xy} \cdot y
  = 2xy + y e^{xy}.
\]

Partial derivative with respect to $y$:
\[
  f_y = \pdv{f}{y}
  = x^2 + e^{xy} \cdot \pdv{}{y}(xy)
  = x^2 + e^{xy} \cdot x
  = x^2 + x e^{xy}.
\]

The total differential is
\[
  df = f_x \, dx + f_y \, dy
  = (2xy + y e^{xy})\,dx + (x^2 + x e^{xy})\,dy.
\]

Explanation: $df$ gives the linear change in $f$ under small changes $dx,dy$.


% 7
\item \textbf{Particle motion on a helix.}

The position is
\[
  \mathbf{r}(t) = a\cos t \,\hat{\imath} + a\sin t \,\hat{\jmath} + ct \,\hat{k},
\]
where $a$ has units of length and $c$ has units of velocity.

\begin{enumerate}
    \item \textit{Type of trajectory.}

    The $x$-$y$ coordinates satisfy
    \[
      x(t)=a\cos t, \qquad y(t)=a\sin t.
    \]
    Then
    \[
      x^2 + y^2 = a^2(\cos^2 t + \sin^2 t)=a^2,
    \]
    which is a circle of radius $a$ in the $xy$-plane. Meanwhile $z(t)=ct$ increases linearly in time. Therefore the path is a helix of radius $a$ rising at constant rate $c$ along the $z$-axis.

    \item \textit{Velocity, acceleration, and speed.}

    Velocity:
    \[
      \mathbf{v}(t) = \dv{\mathbf{r}}{t}
      = -a\sin t \,\hat{\imath} + a\cos t \,\hat{\jmath} + c\,\hat{k}.
    \]

    Acceleration:
    \[
      \mathbf{a}(t) = \dv{\mathbf{v}}{t}
      = -a\cos t \,\hat{\imath} - a\sin t \,\hat{\jmath} + 0\,\hat{k}
      = -a\cos t \,\hat{\imath} - a\sin t \,\hat{\jmath}.
    \]

    Speed (magnitude of velocity):
    \begin{align*}
      v(t) &= \|\mathbf{v}(t)\|
      = \sqrt{(-a\sin t)^2 + (a\cos t)^2 + c^2} \\
      &= \sqrt{a^2(\sin^2 t + \cos^2 t) + c^2} \\
      &= \sqrt{a^2 + c^2}.
    \end{align*}
    The speed is constant in time.

    \item \textit{Directions of $\mathbf{v}$ and $\mathbf{a}$.}

    In the $xy$-plane, $(-a\sin t, a\cos t)$ is tangent to the circle $x^2+y^2=a^2$, and $(-a\cos t, -a\sin t)$ points radially inward toward the center $(0,0)$.

    So:
    \[
      \mathbf{v}(t) \text{ is tangent to the circular motion in the } xy\text{-plane (plus upward motion in } \hat{k}),
    \]
    while
    \[
      \mathbf{a}(t) \text{ points toward the axis of the circle in the } xy\text{-plane (centripetal)}, \text{ with no }\hat{k}\text{ component.}
    \]
    This matches uniform circular motion of radius $a$ in the horizontal projection: velocity is tangent, acceleration is inward.

    \item \textit{Projection of $\mathbf{a}$ onto $\mathbf{v}$. Physical meaning.}

    The projection of $\mathbf{a}$ onto $\mathbf{v}$ is the component of acceleration along the direction of motion. It measures tangential acceleration (change in speed). Compute
    \[
      \text{proj}_{\mathbf{v}}\mathbf{a}
      = \frac{\mathbf{a}\cdot \mathbf{v}}{\|\mathbf{v}\|^2} \,\mathbf{v}.
    \]

    First compute the dot product:
    \begin{align*}
      \mathbf{a}(t)\cdot \mathbf{v}(t)
      &= (-a\cos t)(-a\sin t) + (-a\sin t)(a\cos t) + (0)(c) \\
      &= a^2 \cos t \sin t - a^2 \sin t \cos t \\
      &= 0.
    \end{align*}

    Therefore
    \[
      \text{proj}_{\mathbf{v}}\mathbf{a} = \mathbf{0}.
    \]

    Interpretation: there is no acceleration component along the velocity direction, so the speed is constant. The acceleration is purely normal (centripetal), changing only the direction of the velocity, not its magnitude.
\end{enumerate}


% 8
\item \textbf{Flux through a surface in a vector field.}

Given a flat rectangle in the $xy$-plane with side lengths
\[
a = 0.20~\mathrm{m} \quad \text{(along $x$)}, \qquad
b = 0.30~\mathrm{m} \quad \text{(along $y$)},
\]
and flux density field
\[
\mathbf{F} = 200\,\hat{\imath} + 300\,\hat{\jmath} + 400\,\hat{k}
\quad
\left[\frac{\mathrm{L}}{\mathrm{m}^2\cdot \mathrm{s}}\right].
\]

\begin{enumerate}
    \item \textit{Area vector $\mathbf{A}$.}


    \[
      \mathbf{A} = (a\hat{\imath})\times(b\hat{\jmath}) = (0.20~\mathrm{m})(0.30~\mathrm{m})\hat{k} = (0.060)\hat{k} \quad [\mathrm{m}^2].
    \]

    Its normal vector is in the $+\hat{k}$ direction perpendicular to the surface lying in the $xy$-plane.
    
    \item \textit{Flow rate $Q = \mathbf{F} \cdot \mathbf{A}$.}

    Compute the dot product:
    \begin{align*}
      Q &= \mathbf{F} \cdot \mathbf{A} \\
        &= (200\,\hat{\imath} + 300\,\hat{\jmath} + 400\,\hat{k}) \cdot (0.060\,\hat{k}) \\
        &= 400 \cdot 0.060 \quad \left[\frac{\mathrm{L}}{\mathrm{m}^2\cdot \mathrm{s}}\right] \cdot [\mathrm{m}^2] \\
        &= 24~\mathrm{L/s}.
    \end{align*}

    Units: $\mathbf{F}$ is volume flux per area per time, $\mathbf{A}$ is area, so the product is volume per time.

    \item \textit{Effect of orientation.}

    In general,
    \[
      Q = \mathbf{F} \cdot \mathbf{A}
      = |\mathbf{F}|\,|\mathbf{A}|\cos\theta,
    \]
    where $\theta$ is the angle between $\mathbf{F}$ and the surface normal.
    If you tilt the surface, its area vector $\mathbf{A}$ tilts. The effective flux is reduced by $\cos\theta$.
    \begin{itemize}
        \item If the surface is perpendicular to the flow (normal aligned with $\mathbf{F}$), $\theta=0$, $\cos\theta=1$, and $Q$ is maximal.
        \item If the surface is parallel to the flow (normal perpendicular to $\mathbf{F}$), $\theta=90^\circ$, $\cos\theta=0$, and $Q=0$.
    \end{itemize}
    Here the surface normal is $\hat{k}$, so only the $k$-component of $\mathbf{F}$ contributes to the flux.
\end{enumerate}


% 9
\item \textbf{Velocity, acceleration, and their angle.}

Given
\[
  \mathbf{r}(t) = a t^2 \,\hat{\imath} + b t^3 \,\hat{\jmath}.
\]

\begin{enumerate}
    \item \textit{Velocity and acceleration.}

    Velocity:
    \[
      \mathbf{v}(t) = \dv{\mathbf{r}}{t}
      = 2at \,\hat{\imath} + 3 b t^2 \,\hat{\jmath}.
    \]

    Acceleration:
    \[
      \mathbf{a}(t) = \dv{\mathbf{v}}{t}
      = 2a \,\hat{\imath} + 6 b t \,\hat{\jmath}.
    \]

    \item \textit{Angle between $\mathbf{v}$ and $\mathbf{a}$ at $t=1$.}

    At $t=1$:
    \[
      \mathbf{v}(1) = 2a\,\hat{\imath} + 3b\,\hat{\jmath}, \qquad
      \mathbf{a}(1) = 2a\,\hat{\imath} + 6b\,\hat{\jmath}.
    \]

    The angle $\theta$ between two vectors $\mathbf{u}$ and $\mathbf{w}$ is given by
    \[
      \cos\theta = \frac{\mathbf{u}\cdot\mathbf{w}}{\|\mathbf{u}\|\,\|\mathbf{w}\|}.
    \]

    Compute the dot product:
    \[
      \mathbf{v}(1)\cdot \mathbf{a}(1)
      = (2a)(2a) + (3b)(6b)
      = 4a^2 + 18 b^2.
    \]

    Compute magnitudes:
    \[
      \|\mathbf{v}(1)\| = \sqrt{(2a)^2 + (3b)^2}
      = \sqrt{4a^2 + 9b^2},
    \]
    \[
      \|\mathbf{a}(1)\| = \sqrt{(2a)^2 + (6b)^2}
      = \sqrt{4a^2 + 36b^2}.
    \]

    Therefore
    \[
      \cos\theta
      = \frac{4a^2 + 18b^2}{\sqrt{4a^2 + 9b^2}\,\sqrt{4a^2 + 36b^2}}.
    \]

    So the angle $\theta$ is
    \[
      \theta = \arccos\!\left(
        \frac{4a^2 + 18b^2}{\sqrt{4a^2 + 9b^2}\,\sqrt{4a^2 + 36b^2}}
      \right).
    \]

    Interpretation: the angle depends on $a$ and $b$; if $\mathbf{a}$ is parallel to $\mathbf{v}$, then the direction of motion is speeding up with no change in direction. If not parallel, part of $\mathbf{a}$ turns the velocity.
\end{enumerate}


% 10
\item \textbf{Cross product.}

Given
\[
\mathbf{A} = 2\,\hat{\imath} - 3\,\hat{\jmath} + 1\,\hat{k},
\qquad
\mathbf{B} = -1\,\hat{\imath} + 4\,\hat{\jmath} + 2\,\hat{k}.
\]


The cross product $\mathbf{A}\times\mathbf{B}$ is


\[
\begin{aligned}
\mathbf{A}\times\mathbf{B}
&= (2\hat{\imath}-3\hat{\jmath}+\hat{k})\times(-\hat{\imath}+4\hat{\jmath}+2\hat{k}) \\
&= 2\hat{\imath}\times(-\hat{\imath}) + 2\hat{\imath}\times 4\hat{\jmath} + 2\hat{\imath}\times 2\hat{k}
  -3\hat{\jmath}\times(-\hat{\imath}) -3\hat{\jmath}\times 4\hat{\jmath} -3\hat{\jmath}\times 2\hat{k} \\
&\quad + \hat{k}\times(-\hat{\imath}) + \hat{k}\times 4\hat{\jmath} + \hat{k}\times 2\hat{k} \\
&= 0 + 8\hat{k} - 4\hat{\jmath} - 3\hat{k} + 0 - 6\hat{\imath} - \hat{\jmath} - 4\hat{\imath} + 0 \\
&= -10\,\hat{\imath} - 5\,\hat{\jmath} + 5\,\hat{k}.
\end{aligned}
\]

Explanation: $\mathbf{A}\times\mathbf{B}$ is perpendicular to both $\mathbf{A}$ and $\mathbf{B}$ and its magnitude is $|\mathbf{A}||\mathbf{B}|\sin\phi$, where $\phi$ is the angle between them.
\end{enumerate}


\end{document}
